
We presented a methodological framework to test whether multi-objective code modifications exhibit aspect-dominant structure that would justify architectural factorization. Through rigorous validation combining spectral analysis, permutation testing, and cross-validation, we provide methodology to answer this question decisively.

Our proof-of-concept demonstration on synthetic test data validates that the framework works as designed—correctly identifying independence without factorization when this pattern exists. The synthetic data exhibited statistical independence (coupling = 0.072) yet no aspect-aligned directional structure (spectral gap = 1.580 < 2.0, permutation p = 0.955). This demonstrates the key conceptual distinction: independence does not imply factorization.

The methodological contribution is establishing how to test empirical separability before architectural commitment. Multi-task learning theory suggests independent tasks should have separable subspaces, justifying architectural constraints. Our framework provides tools to validate this assumption for code rather than borrowing intuitions from other domains.

We acknowledge the critical limitation: our evaluation used synthetic test data for methodological validation, not real GitHub commits. The methodological rigor—residual covariance analysis, permutation testing with 1000 shuffles, directional stability measures, cross-validation—provides a validated template ready for real data application. Real data collection is the required next step to determine whether independence without factorization characterizes actual developer behavior.

Beyond the methodological framework, we propose four constructive alternative structural hypotheses: hierarchical quality (DAG dependencies), contextual factorization (domain-specific clusters), commit size dependence (scale-dependent structure), and temporal dynamics (sequential patterns). These represent genuine research directions that would inform architectural decisions if validated.

The path forward is clear. For practitioners: await empirical validation before committing to architectural complexity. For researchers: apply this framework to real data across languages, domains, and commit scales. Test alternative structural hypotheses rigorously. Let data guide architectural decisions rather than borrowed intuitions.

The broader lesson: assumptions transferred from adjacent domains require empirical validation. The assumption that task independence implies architectural factorization may hold in vision/NLP but requires testing for code. Methodological rigor accelerates scientific progress by revealing when complexity is justified and when simplicity suffices.
